% inciso-a.tex
% Fragmento LaTeX. Pensado para incluirse vía % inciso-a.tex
% Fragmento LaTeX. Pensado para incluirse vía % inciso-a.tex
% Fragmento LaTeX. Pensado para incluirse vía % inciso-a.tex
% Fragmento LaTeX. Pensado para incluirse vía \input{inciso-a} en un documento principal.
% Requiere: codificación UTF-8 (\usepackage[utf8]{inputenc}) y soporte de español
% (\usepackage[spanish]{babel} o polyglossia).

En una arquitectura dirigida por eventos con topología de \emph{Broker}, el flujo de
mensajes se distribuye entre procesadores independientes a través de un intermediario
ligero, sin que exista un componente central que orqueste la transacción global.
Cuando un procesador falla de forma repetida sobre un mensaje concreto debido a un
error en los datos, se desencadenan varios efectos en cadena que conviene analizar
por separado.

\subsubsection{Invisibilidad del fallo}

El primero y más problemático es la \textbf{invisibilidad del fallo}. Al no haber un
mediador con visión global, ningún otro componente del sistema se entera de que el
procesador ha caído ni de que el flujo de negocio asociado a ese mensaje se ha
detenido. La transacción queda atrapada en un estado incompleto sin que se dispare
ninguna alerta automática. Este bloqueo silencioso es precisamente la consecuencia
directa de la ausencia de orquestador.

\subsubsection{Problemas de consistencia}

A esto se suma un problema de \textbf{consistencia}. El resto de procesadores
continúa ejecutándose con normalidad ---registrando pagos, modificando inventarios,
emitiendo notificaciones paralelas--- sin saber que una fase crítica ha fallado. Se
producen así asimetrías y estados inconsistentes en todo el ecosistema distribuido,
lo que agrava el modelo de consistencia eventual inherente a la arquitectura. Cabe
recordar que en esta topología el \emph{broker} es ligero y carece de capacidades de
coordinación transaccional (\emph{two-phase commit}), por lo que tampoco es posible
garantizar atomicidad entre varios procesadores.

\subsubsection{Impacto en el usuario final}

Desde el punto de vista del usuario final, la experiencia se degrada de forma
notable. Al operar de manera totalmente asíncrona y sin canal de retorno directo, el
usuario no recibe ninguna retroalimentación sobre el fallo de su solicitud: la
transacción queda en una especie de limbo indefinido, lo que afecta a la confianza
y a la percepción general del servicio.

\subsubsection{Efecto poison pill y degradación del sistema}

Por último, el mensaje que origina el fallo se comporta como un clásico
\emph{poison pill}: un mensaje que nunca podrá procesarse correctamente y que, si el
procesador intenta reintentos síncronos o bloqueantes, acaba deteniendo la lectura
de la cola (\emph{head-of-line blocking}). Esto degrada drásticamente la
\emph{responsiveness} del sistema para los mensajes posteriores y, además, consume
recursos ---CPU, memoria, conexiones--- en cada reintento inútil, con riesgo de
entrar en un bucle infinito si no se define una política de máximo de reintentos.
 en un documento principal.
% Requiere: codificación UTF-8 (\usepackage[utf8]{inputenc}) y soporte de español
% (\usepackage[spanish]{babel} o polyglossia).

En una arquitectura dirigida por eventos con topología de \emph{Broker}, el flujo de
mensajes se distribuye entre procesadores independientes a través de un intermediario
ligero, sin que exista un componente central que orqueste la transacción global.
Cuando un procesador falla de forma repetida sobre un mensaje concreto debido a un
error en los datos, se desencadenan varios efectos en cadena que conviene analizar
por separado.

\subsubsection{Invisibilidad del fallo}

El primero y más problemático es la \textbf{invisibilidad del fallo}. Al no haber un
mediador con visión global, ningún otro componente del sistema se entera de que el
procesador ha caído ni de que el flujo de negocio asociado a ese mensaje se ha
detenido. La transacción queda atrapada en un estado incompleto sin que se dispare
ninguna alerta automática. Este bloqueo silencioso es precisamente la consecuencia
directa de la ausencia de orquestador.

\subsubsection{Problemas de consistencia}

A esto se suma un problema de \textbf{consistencia}. El resto de procesadores
continúa ejecutándose con normalidad ---registrando pagos, modificando inventarios,
emitiendo notificaciones paralelas--- sin saber que una fase crítica ha fallado. Se
producen así asimetrías y estados inconsistentes en todo el ecosistema distribuido,
lo que agrava el modelo de consistencia eventual inherente a la arquitectura. Cabe
recordar que en esta topología el \emph{broker} es ligero y carece de capacidades de
coordinación transaccional (\emph{two-phase commit}), por lo que tampoco es posible
garantizar atomicidad entre varios procesadores.

\subsubsection{Impacto en el usuario final}

Desde el punto de vista del usuario final, la experiencia se degrada de forma
notable. Al operar de manera totalmente asíncrona y sin canal de retorno directo, el
usuario no recibe ninguna retroalimentación sobre el fallo de su solicitud: la
transacción queda en una especie de limbo indefinido, lo que afecta a la confianza
y a la percepción general del servicio.

\subsubsection{Efecto poison pill y degradación del sistema}

Por último, el mensaje que origina el fallo se comporta como un clásico
\emph{poison pill}: un mensaje que nunca podrá procesarse correctamente y que, si el
procesador intenta reintentos síncronos o bloqueantes, acaba deteniendo la lectura
de la cola (\emph{head-of-line blocking}). Esto degrada drásticamente la
\emph{responsiveness} del sistema para los mensajes posteriores y, además, consume
recursos ---CPU, memoria, conexiones--- en cada reintento inútil, con riesgo de
entrar en un bucle infinito si no se define una política de máximo de reintentos.
 en un documento principal.
% Requiere: codificación UTF-8 (\usepackage[utf8]{inputenc}) y soporte de español
% (\usepackage[spanish]{babel} o polyglossia).

En una arquitectura dirigida por eventos con topología de \emph{Broker}, el flujo de
mensajes se distribuye entre procesadores independientes a través de un intermediario
ligero, sin que exista un componente central que orqueste la transacción global.
Cuando un procesador falla de forma repetida sobre un mensaje concreto debido a un
error en los datos, se desencadenan varios efectos en cadena que conviene analizar
por separado.

\subsubsection{Invisibilidad del fallo}

El primero y más problemático es la \textbf{invisibilidad del fallo}. Al no haber un
mediador con visión global, ningún otro componente del sistema se entera de que el
procesador ha caído ni de que el flujo de negocio asociado a ese mensaje se ha
detenido. La transacción queda atrapada en un estado incompleto sin que se dispare
ninguna alerta automática. Este bloqueo silencioso es precisamente la consecuencia
directa de la ausencia de orquestador.

\subsubsection{Problemas de consistencia}

A esto se suma un problema de \textbf{consistencia}. El resto de procesadores
continúa ejecutándose con normalidad ---registrando pagos, modificando inventarios,
emitiendo notificaciones paralelas--- sin saber que una fase crítica ha fallado. Se
producen así asimetrías y estados inconsistentes en todo el ecosistema distribuido,
lo que agrava el modelo de consistencia eventual inherente a la arquitectura. Cabe
recordar que en esta topología el \emph{broker} es ligero y carece de capacidades de
coordinación transaccional (\emph{two-phase commit}), por lo que tampoco es posible
garantizar atomicidad entre varios procesadores.

\subsubsection{Impacto en el usuario final}

Desde el punto de vista del usuario final, la experiencia se degrada de forma
notable. Al operar de manera totalmente asíncrona y sin canal de retorno directo, el
usuario no recibe ninguna retroalimentación sobre el fallo de su solicitud: la
transacción queda en una especie de limbo indefinido, lo que afecta a la confianza
y a la percepción general del servicio.

\subsubsection{Efecto poison pill y degradación del sistema}

Por último, el mensaje que origina el fallo se comporta como un clásico
\emph{poison pill}: un mensaje que nunca podrá procesarse correctamente y que, si el
procesador intenta reintentos síncronos o bloqueantes, acaba deteniendo la lectura
de la cola (\emph{head-of-line blocking}). Esto degrada drásticamente la
\emph{responsiveness} del sistema para los mensajes posteriores y, además, consume
recursos ---CPU, memoria, conexiones--- en cada reintento inútil, con riesgo de
entrar en un bucle infinito si no se define una política de máximo de reintentos.
 en un documento principal.
% Requiere: codificación UTF-8 (\usepackage[utf8]{inputenc}) y soporte de español
% (\usepackage[spanish]{babel} o polyglossia).

En una arquitectura dirigida por eventos con topología de \emph{Broker}, el flujo de
mensajes se distribuye entre procesadores independientes a través de un intermediario
ligero, sin que exista un componente central que orqueste la transacción global.
Cuando un procesador falla de forma repetida sobre un mensaje concreto debido a un
error en los datos, se desencadenan varios efectos en cadena que conviene analizar
por separado.

\subsubsection{Invisibilidad del fallo}

El primero y más problemático es la \textbf{invisibilidad del fallo}. Al no haber un
mediador con visión global, ningún otro componente del sistema se entera de que el
procesador ha caído ni de que el flujo de negocio asociado a ese mensaje se ha
detenido. La transacción queda atrapada en un estado incompleto sin que se dispare
ninguna alerta automática. Este bloqueo silencioso es precisamente la consecuencia
directa de la ausencia de orquestador.

\subsubsection{Problemas de consistencia}

A esto se suma un problema de \textbf{consistencia}. El resto de procesadores
continúa ejecutándose con normalidad ---registrando pagos, modificando inventarios,
emitiendo notificaciones paralelas--- sin saber que una fase crítica ha fallado. Se
producen así asimetrías y estados inconsistentes en todo el ecosistema distribuido,
lo que agrava el modelo de consistencia eventual inherente a la arquitectura. Cabe
recordar que en esta topología el \emph{broker} es ligero y carece de capacidades de
coordinación transaccional (\emph{two-phase commit}), por lo que tampoco es posible
garantizar atomicidad entre varios procesadores.

\subsubsection{Impacto en el usuario final}

Desde el punto de vista del usuario final, la experiencia se degrada de forma
notable. Al operar de manera totalmente asíncrona y sin canal de retorno directo, el
usuario no recibe ninguna retroalimentación sobre el fallo de su solicitud: la
transacción queda en una especie de limbo indefinido, lo que afecta a la confianza
y a la percepción general del servicio.

\subsubsection{Efecto poison pill y degradación del sistema}

Por último, el mensaje que origina el fallo se comporta como un clásico
\emph{poison pill}: un mensaje que nunca podrá procesarse correctamente y que, si el
procesador intenta reintentos síncronos o bloqueantes, acaba deteniendo la lectura
de la cola (\emph{head-of-line blocking}). Esto degrada drásticamente la
\emph{responsiveness} del sistema para los mensajes posteriores y, además, consume
recursos ---CPU, memoria, conexiones--- en cada reintento inútil, con riesgo de
entrar en un bucle infinito si no se define una política de máximo de reintentos.
